%% kapa: exam=1 round=33 year=2022 period=2 form=A subject=accounting subject_ko=회계학 seq=21 number=61
%% source: Q-Net 감정평가사 기출문제 (cst00309) artlSeq=5209997 file=raw/1008/5209997/2022년도 제33회 감정평가사 1차시험 2교시(감정평가 관계법규, 회계학).pdf
%% conversion: kordoc 4.24.51 -> tools/gen1.py (auto); review=proofread-sonnet; flags=none
\begin{question}{61}
(주)감평은 특정차입금 없이 일반차입금을 사용하여 건물을 신축하였다. 건물은 차입원가 자본화 대상인 적격자산이다. 신축 건물과 관련한 자료가 다음과 같을 경우, 20×1년도에 자본화할 차입원가(A)와 20×2년도에 자본화할 차입원가(B)는? (단, 계산시 월할 계산하며, 전기에 자본화한 차입원가는 적격자산의 연평균 지출액 계산 시 포함하지 않는다.)
\begin{framedtext}
○ 공사기간: 20×1년 5월 1일 \textasciitilde{} 20×2년 6월 31일 % [원문 그대로]
\par
○ 공사비 지출: \par
{\small\centering\begin{tabular}{ccc}
20×1년 5월 1일 & 20×1년 10월 1일 & 20×2년 4월 1일 \\
￦300,000 & ￦200,000 & ￦100,000 \\
\end{tabular}\par}
○ 일반차입금 자본화 연이자율: \par
{\small\centering\begin{tabular}{cc}
20×1년 & 20×2년 \\
10\% & 8\% \\
\end{tabular}\par}
○ 실제 발생한 이자비용: \par
{\small\centering\begin{tabular}{cc}
20×1년 & 20×2년 \\
￦20,000 & ￦24,200 \\
\end{tabular}\par}
\end{framedtext}
\begin{choices}
  \item (A) ￦20,000 \quad (B) ￦22,000
  \item (A) ￦20,000 \quad (B) ￦24,200
  \item (A) ￦20,000 \quad (B) ￦25,000
  \item (A) ￦25,000 \quad (B) ￦22,000
  \item (A) ￦25,000 \quad (B) ￦24,200
\end{choices}
\end{question}
\answer{1}
